\abstract{РЕФЕРАТ}

Объем работы равен \formbytotal{lastpage}{страниц}{е}{ам}{ам}. Работа содержит \formbytotal{figurecnt}{иллюстраци}{ю}{и}{й}, \formbytotal{tablecnt}{таблиц}{у}{ы}{}, \arabic{bibcount} библиографических источников и \formbytotal{числоПлакатов}{лист}{}{а}{ов} графического материала. Количество приложений – 2. Графический материал представлен в приложении А. Фрагменты исходного кода представлены в приложении Б.

Перечень ключевых слов: игровой движок, OpenGL, C\#, Windows Forms, рендеринг, шейдеры, трехмерная графика, редактор сцен, импорт моделей, STL, GLB, сериализация сцены, система освещения, модель Блинна-Фонга, кватернионы, рейкастинг, игровые объекты, трансформация, камера.

Объектом разработки является программный комплекс с визуальной средой проектирования для создания компьютерных игр, включающий игровой движок на языке C\# с использованием графического API OpenGL и редактор сцен на платформе Windows Forms.

Целью выпускной квалификационной работы является разработка и внедрение игрового движка для создания 3D интерактивных приложений в реальном времени.

В процессе создания программного комплекса были реализованы основные компоненты игрового движка: графическая подсистема, система управления игровыми объектами, менеджер сцен, система камер и освещения. Разработана подсистема импорта  моделей в форматах STL и GLB с поддержкой материалов и текстур. Реализован редактор сцен, обеспечивающий визуальное управление игровыми объектами, настройку их трансформации, тегирование, а также импорт моделей и сохранение сцен.

При разработке программного комплекса использовались язык программирования C\#, графический API OpenGL, библиотеки OpenTK, SixLabors.ImageSharp, Windows Forms, а также инструменты сборки .NET 6.0 и выше.

На разработанном движке была успешно создана тестовая 3D игра с механикой выбора и возрождения объектов.

\selectlanguage{english}
\abstract{ABSTRACT}
  
The volume of work is \formbytotal{lastpage}{page}{}{s}{s}. The work contains \formbytotal{figurecnt}{illustration}{}{s}{s}, \formbytotal{tablecnt}{table}{}{s}{s}, \arabic{bibcount} bibliographic sources and \formbytotal{числоПлакатов}{sheet}{}{s}{s} of graphic material. The number of applications is 2. The graphic material is presented in annex A. The layout of the site, including the connection of components, is presented in annex B.

Keywords: game engine, OpenGL, C\#, Windows Forms, rendering, shaders, three-dimensional graphics, scene editor, model import, STL, GLB, scene serialization, lighting system, Blinn-Fong model, quaternions, raycasting, game objects, transformation, camera.

The object of development is a software package with a visual design environment for creating computer games, including a game engine in C using the OpenGL graphics API and a scene editor on the Windows Forms platform.

The purpose of the final qualification is to develop and implement a game engine for creating 3D interactive applications in real time.

During the creation of the software package, the main components of the game engine were implemented: a graphics subsystem, a game object management system, a scene manager, a camera and lighting system. A subsystem for importing models in STL and GLB formats with support for materials and textures has been developed. A scene editor has been implemented that provides visual control of game objects, setting up their transformations, tagging, as well as importing models and saving scenes.

When developing the software package, the C\# programming language, the OpenGL graphics API, the OpenTK libraries, SixLabors.ImageSharp, Windows Forms, and assembly tools were used.NET 6.0 and higher.

A 3D test game with the mechanics of selecting and reviving objects was successfully created using the developed engine.
\selectlanguage{russian}
